% !TeX root = 流体力学笔记.tex
\begin{titlepage}
	\begin{tikzpicture}[remember picture, overlay]
		% ---- 1. 深蓝渐变铺满整个封面（海洋深蓝色调）----
		\shade[top color=DarkBlue, middle color=DarkBlue!90, bottom color=DarkBlue!65!black]
			(current page.south west) rectangle (current page.north east);

		% ---- 2. 细点阵（模拟网格，微弱纹理）----
		\foreach \x in {1,...,20} {
			\foreach \y in {1,...,29} {
				\fill[white, opacity=0.05] (\x,\y) circle[radius=0.018cm];
			}
		}

		% ---- 3. 主标题后方的柔光晕 ----
		\fill[white, opacity=0.05] ([yshift=-8.2cm]current page.north) circle[radius=7.0cm];
		\fill[white, opacity=0.05] ([yshift=-8.2cm]current page.north) circle[radius=4.6cm];

		% ---- 4. 右上：理想流体绕圆柱流动流线图 ----
		\begin{scope}[shift={([xshift=-5.8cm,yshift=-4.8cm]current page.north east)}]
			% 圆柱截面
			\fill[white, opacity=0.15] (0,0) circle[radius=0.8cm];
			\draw[white, line width=1.0pt, opacity=0.8] (0,0) circle[radius=0.8cm];
			% 流线（水平绕流）
			\foreach \y in {-2.0,-1.5,-1.0,-0.5,0.5,1.0,1.5,2.0} {
				\draw[white, opacity=0.6, line width=0.7pt, domain=-3:3, samples=100, smooth]
					plot (\x, {\y + 0.8*\y/(\x*\x+\y*\y+0.01)});
			}
			% 中心流线（绕过圆柱）
			\draw[white, opacity=0.7, line width=0.9pt, domain=-3:-0.9, samples=60, smooth]
				plot (\x, {0.8*0.1/(\x*\x+0.01)});
			\draw[white, opacity=0.7, line width=0.9pt, domain=0.9:3, samples=60, smooth]
				plot (\x, {0.8*0.1/(\x*\x+0.01)});
			% 标签
			\node[anchor=south west, text=white, opacity=0.7, font=\footnotesize\itshape] at (3.0,2.2) {流线};
		\end{scope}

		% ---- 5. 左下：纳维-斯托克斯方程 ----
		\begin{scope}[shift={([xshift=1.0cm,yshift=1.0cm]current page.south west)}]
			\node[anchor=north west, text=white, opacity=0.85, align=left] at (0,2.5) {
				{\small\itshape$\rho\left(\frac{\partial \bm{v}}{\partial t}+\bm{v}\cdot\nabla\bm{v}\right)$}\\[0.3em]
				{\small\itshape$\quad =-\nabla p+\mu\nabla^{2}\bm{v}+\bm{f}$}
			};
			\node[anchor=north west, text=white, opacity=0.6, font=\footnotesize] at (0,0) {Navier-Stokes};
		\end{scope}

		% ---- 6. 金色细边框 ----
		\draw[gold, line width=0.7pt, opacity=0.55]
			([shift={(0.85,0.85)}]current page.south west) rectangle
			([shift={(-0.85,-0.85)}]current page.north east);

		% ---- 7. 顶部眉题（英文小字 + 金色短线）----
		\coordinate (T) at ([yshift=-3.5cm]current page.north);
		\draw[gold, line width=1.0pt] (T) ++(-3.4cm,0) -- ++(1.0cm,0);
		\draw[gold, line width=1.0pt] (T) ++(2.4cm,0) -- ++(1.0cm,0);
		\node[anchor=north, text=white, opacity=0.85, align=center] at (T) {
			{\sffamily\fontsize{13}{16}\selectfont
				\uppercase{Fluid \ Mechanics}}
		};

		% ---- 8. 主标题（居中，使用 \notename）----
		\node[anchor=center, align=center, text=white] at ([yshift=-8.0cm]current page.north) {%
			{\fontsize{40}{50}\sffamily\bfseries \notename}
		};

		% ---- 9. 金色菱形 + 双细线 ----
		\coordinate (C) at ([yshift=-11.4cm]current page.north);
		\draw[gold, line width=1.1pt] (C) ++(-2.2cm,0) -- ++(1.55cm,0);
		\draw[gold, line width=1.1pt] (C) ++(0.65cm,0) -- ++(1.55cm,0);
		\node[fill=gold, minimum size=0.24cm, inner sep=0pt, rotate=45] at (C) {};

		% ---- 10. 副标题 ----
		\node[anchor=north, text=white, opacity=0.9, align=center]
			at ([yshift=-12.2cm]current page.north) {
			{\fontsize{16}{20}\sffamily 学\ \ 习\ \ 笔\ \ 记}
		};

		% ---- 11. 底部作者、日期信息 ----
		\coordinate (B) at ([yshift=3.4cm]current page.south);
		\draw[gold, line width=0.8pt, opacity=0.7] (B) ++(-1.1cm,0.85cm) -- ++(2.2cm,0);
		\node[anchor=south, text=white, align=center] at (B) {
			\begin{tabular}{c}
				{\fontsize{20}{24}\sffamily\bfseries \noteauthor} \\[0.4cm]
				{\large \sffamily \today}
			\end{tabular}
		};
	\end{tikzpicture}
\end{titlepage}
